\documentclass[10pt,a4paper]{amsart}
\usepackage[margin=19mm]{geometry}
\usepackage{amsmath,amssymb,mathtools}
\usepackage[scr=rsfs,cal=euler]{mathalfa}
\usepackage{xfrac}
\usepackage[unicode,hidelinks,bookmarks=false]{hyperref}
\newcommand{\ie}{i.e.\ }
\newcommand{\cf}{cf.\ }
\newcommand{\resp}{respectively\ }
\newcommand{\Subsection}{\subsection}
\providecommand{\nobreakdash}{\nobreak\hbox{-}}
\setlength{\emergencystretch}{2em}
\multlinegap0pt
\usepackage{amssymb}
\usepackage{mathtools}
\usepackage[shortlabels]{enumitem}
\usepackage{tikz}
\usepackage{float}
\newtheorem{lemma}{Lemma}
\title{Minimal Equiangular Hyperbolic Polyhedra\\ in the Tetrahedral Range}
\author{Andrey Egorov}
\date{Source: arXiv:2608.14223v2 (CC BY 4.0)}
\begin{document}
\maketitle

\noindent\emph{Source and licence.} Minimal Equiangular Hyperbolic Polyhedra\\ in the Tetrahedral Range. Version arXiv:2608.14223v2 (CC BY 4.0), distributed by arXiv under the Creative Commons Attribution 4.0 International licence, http://creativecommons.org/licenses/by/4.0/, as recorded in the arXiv licence metadata for this exact version and shown on the abstract page https://arxiv.org/abs/2608.14223v2. Full licence notice: \texttt{licenses/arxiv-2608.14223-CC-BY-4.0.txt}.\par\smallskip

\noindent\textit{Editorial scope.} Counted unit: the article's lemma lem:reflex-patch-crossing (source0.tex lines 598--632). The article's complete original proof of it is delivered with it (source0.tex lines 634--798, one \texttt{proof} environment of 165 lines). No mathematics is written, repaired or re-derived: the statement and the proof are the article's own lines, in the article's own notation, and no retained line is substituted or re-flowed.  No further source-local context is needed: every cross-reference of every retained line resolves inside the two delivered spans.  Nothing was omitted from any retained span, so the package carries no omission record.  No retained line contains a citation, so the wrapper carries no bibliography and no bibliography entry.  No retained line contains a figure environment, an image-inclusion command or drawing code, so no image is delivered and the package declares no asset dependency.  Editorial formatting only, in addition: an AMS \texttt{amsart} preamble that restates the article's own declarations and macros used by the retained lines, namely the environments \texttt{lemma} on separate counters, together with the standard packages those lines need.  The retained environments print in this document as Lemma 1 in order of appearance, whereas the article prints the same items with the numbering of its own full document.  The complete native source of the article as distributed in the arXiv e-print for this exact version is bundled unchanged as \texttt{sources/207647/source0.tex}.
\begin{lemma}[Crossing the reflex two-triangle patch]
\label{lem:reflex-patch-crossing}
Let \(0<s\le L\), and let
\[
        K_s=ABC\cup_{AB}ABD
\]
be the union of two spherical triangles with
\[
 |AB|=s,\qquad |AC|=|BC|=|AD|=|BD|=L,
\]
developed on opposite sides of \(AB\).  If a transverse local geodesic segment
\(\kappa\) crosses, in this order, the sides
\[
        AC,\quad AB,\quad AD,
\]
then
\[
        \ell(\kappa)>\pi.
\]
The symmetric assertion with \(BC,BA,BD\) in place of \(AC,AB,AD\) also
holds.

The same strict estimate holds for every nonconstant one-sided limit of either
of these transverse passages.  In particular, the marked intersections with
\(AC,AB,AD\), or with \(BC,BA,BD\), may converge to endpoints of those sides,
and the limiting passage may run along a boundary edge of its carrier.

There is also a vertex version. Suppose that a local geodesic in \(K_s\)
passes through \(A\), with one germ in \(ABC\) and the other in \(ABD\), and
that the link angle between these germs on the \(K_s\)-side is at least
\(\pi\). Its
subarc from \(BC\) to \(BD\) through \(A\) has length \(>\pi\).
Symmetrically, a local geodesic through \(B\) from \(AC\) to \(AD\) has
length \(>\pi\).
\end{lemma}

\begin{proof}
Put
\[
 q=\sin\frac{s}{2},\qquad r=\cos\frac{s}{2},\qquad
 z=\sqrt{1-c^2-q^2},
\]
where \(c=\cos L<0\).  Since
\[
 q^2\le\sin^2\frac L2=\frac{1-c}{2},
\]
we have
\[
 z^2\ge\frac{(1-c)(1+2c)}2>0.
\]
After a rotation of \(S^2\), the developed vertices are
\[
 \begin{aligned}
 A&=(r,-q,0),& B&=(r,q,0),\\
 C&=(c/r,0,z/r),&D&=(c/r,0,-z/r).
 \end{aligned}
\]
Indeed, these vectors are unit vectors,
\[
 A\cdot B=\cos s,
 \qquad
 A\cdot C=B\cdot C=A\cdot D=B\cdot D=c,
\]
and \(D\) is the reflection of \(C\) across the plane of the great circle
containing \(AB\).  Thus the coordinates realize exactly the two prescribed
triangles.

They also give the positive linear relation
\begin{equation}\label{eq:reflex-patch-linear-relation}
        A+B+\mu(C+D)=0,
        \qquad
        \mu=-\frac{r^2}{c}>0 .
\end{equation}

Let \(P\in AC\), \(Q\in AB\), and \(R\in AD\) be the three crossing points.
The developed segment lies on a great circle
\(G=\nu^\perp\cap S^2\).  Choosing the sign of its normal \(\nu\), write
\[
 \nu\cdot A=-a<0,\qquad
 \nu\cdot B=b>0,\qquad
 \nu\cdot C=u>0,\qquad
 \nu\cdot D=v>0.
\]
The signs follow from the three transverse crossings.  Unnormalized vectors
on the rays through the crossing points are
\[
 p_0=uA+aC,\qquad
 q_0=bA+aB,\qquad
 r_0=vA+aD.
\]
All coefficients are positive, so their normalizations are precisely the
points \(P,Q,R\) on the short spherical sides.

Taking the scalar product of
\eqref{eq:reflex-patch-linear-relation} with \(\nu\) gives
\[
        b-a=-\mu(u+v).
\]
Using \(B=-A-\mu(C+D)\), we therefore obtain
\[
 \begin{aligned}
 q_0
 &=bA+aB\\
 &=(b-a)A-a\mu(C+D)\\
 &=-\mu\bigl((u+v)A+a(C+D)\bigr)\\
 &=-\mu(p_0+r_0).
 \end{aligned}
\]
The points \(P\) and \(R\) cannot be antipodal.  Otherwise their common line
would belong to both planes \(\operatorname{span}(A,C)\) and
\(\operatorname{span}(A,D)\), whose intersection is
\(\operatorname{span}(A)\) because the displayed coordinates make
\(A,C,D\) linearly independent.  This would force \(P=\pm A\), contrary to
transversality.

Consequently the antipodal point \(-Q\) lies in the interior of the short
spherical arc from \(P\) to \(R\): its direction is a positive linear
combination of the directions of \(P\) and \(R\).  Hence \(Q\) lies on the
complementary long arc from \(P\) to \(R\).  Since the developed segment
\(\kappa\) passes through \(Q\), it is this long arc, and
\[
        \ell(\kappa)=2\pi-d_{S^2}(P,R)>\pi.
\]
The case on the \(B\)-side follows by interchanging \(A\) and \(B\).

We next justify the assertion about one-sided limits.  Let
\(\kappa_n\) be transverse passages of the type just considered, with marked
points \(P_n\in AC\), \(Q_n\in AB\), and \(R_n\in AD\), and suppose that they
converge from the chosen side to a nonconstant limiting passage \(\kappa\).
The transverse part of the proof shows that \(\kappa_n\) is the long arc from
\(P_n\) to \(R_n\); hence
\[
        \ell(\kappa_n)=2\pi-d_{S^2}(P_n,R_n).
\]
After passing to the limit, write \(P_n\to P\in AC\) and
\(R_n\to R\in AD\).  The points \(P\) and \(R\) cannot be distinct and
antipodal.  Indeed, their common line would then lie in both
\(\operatorname{span}(A,C)\) and \(\operatorname{span}(A,D)\).  As above, the
intersection of these two planes is \(\operatorname{span}(A)\), so one of the
two points would be \(-A\); but \(-A\) lies on neither of the short sides
\(AC,AD\).  Therefore \(d_{S^2}(P,R)<\pi\).  The choice of one side fixes the
long, rather than the short, limiting arc, and continuity gives
\[
        \ell(\kappa)=2\pi-d_{S^2}(P,R)>\pi.
\]
If \(P=R\), the limiting oriented arc is a full great circle and has length
\(2\pi\), so the same conclusion holds.  This argument also covers a marked
point arriving at a vertex or a limiting great circle coinciding with a
boundary side: only the endpoints and the chosen long arc enter the length
calculation.  The \(B\)-side limits are identical by symmetry.

It remains to prove the vertex version. Let \(P\in BC\), \(R\in BD\), and
put
\[
        x=d(A,P),\qquad y=d(A,R).
\]
Let \(u=\angle PAB\) and \(v=\angle BAR\), measured inside the two special
triangles.  Then
\[
        0\le u,v\le\delta<\pi,
        \qquad
        \pi\le u+v<2\pi,
\]
where the lower bound for the sum is local geodesicity at \(A\).

If \(s\ge\pi/2\), every point of \(BC\) and \(BD\) is at distance at least
\(\pi/2\) from \(A\). Indeed, its representing vector is the normalization
of a nonnegative linear combination of the endpoint vectors, and their scalar
products with \(A\) are \(\cos s\le0\) and \(c<0\). Equality can occur only
when \(s=\pi/2\) and the point is \(B\). It cannot occur for both \(P\) and
\(R\), since then \(u=v=0\), contrary to \(u+v\ge\pi\). Hence
\(x+y>\pi\).

Suppose that \(s<\pi/2\). In the tangent plane at \(A\), take the direction
of \(AB\) as angle zero. Intersecting the great-circle ray which makes angle
\(u\) with \(AB\) with the great circle \(BC\), using the displayed
coordinates, gives
\[
        \cot x=\cot s\cos u+\frac{c}{2rz}\sin u.
\]
For completeness, write the ray as
\(\cos x\,A+\sin x\,\tau(u)\), take its scalar product with \(B\times C\),
and solve the resulting linear equation for \(\cot x\).
The identical calculation in \(ABD\) gives
\[
        \cot y=\cot s\cos v+\frac{c}{2rz}\sin v.
\]
Here \(\cot s>0\), \(c/(2rz)<0\), and
\[
 \cos u+\cos v
 =2\cos\frac{u+v}{2}\cos\frac{u-v}{2}\le0,
 \qquad
 \sin u+\sin v>0.
\]
Indeed, \(|u-v|<\pi\), so the second cosine in the product is positive.
Consequently \(\cot x+\cot y<0\). Since \(0<x,y<\pi\),
\[
 \sin(x+y)=\sin x\sin y\,(\cot x+\cot y)<0,
\]
and therefore \(x+y>\pi\). The assertion through \(B\) follows by symmetry.
\end{proof}

\end{document}
